%! Graphing Radical Functions \& Transformations — Unit 5, Lesson 5.4 — Exit Ticket (Answer Key)
\documentclass[10pt]{article}
\usepackage{algebra2-article}
\usepackage{algebra2-key}   % pulls in -boxes; do NOT also load -boxes
\newcommand{\TallMath}[1]{$\displaystyle #1\rule[-1.4em]{0pt}{3.2em}$}
% In-formula answer slot (\ans is text-mode and must never appear inside $...$).
\newcommand{\mans}[1]{{\color{keyred}\mathbf{#1}}}
\newcommand{\lgrid}[4]{%
  \draw[step=1, linegray!40, very thin] (#1,#3) grid (#2,#4);
  \draw[->, semithick, charcoal] (#1,0)--(#2,0) node[below right, font=\scriptsize]{$x$};
  \draw[->, semithick, charcoal] (0,#3)--(0,#4) node[left, font=\scriptsize]{$y$};}

\begin{document}

\pageheader{Unit 5, Lesson 5.4}{Exit Ticket --- Answer Key}
\namedateperiod

\vspace{0.08in}

No notes. Write every domain and range in \textbf{interval notation}.

\begin{enumerate}[leftmargin=1.8em, itemsep=9pt]
  \item \textbf{From the equation.} \; $f(x)=-3\sqrt{x+2}+4$

        \vspace{3pt}
        (a) Anchor point: \ans{$(-2,4)$} \hspace{0.3in}
        (b) Domain: \ans{$[-2,\infty)$} \hspace{0.3in}
        (c) Range: \ans{$(-\infty,4]$}

        \vspace{5pt}
        (d) The anchor is an absolute \ans{maximum}, and I knew that from \ans{the negative $a=-3$}.

        \vspace{5pt}
        (e) Describe all three transformations from the parent $y=\sqrt{x}$. \ansline{Left $2$; a
        vertical stretch by a factor of $3$ combined with a reflection across the $x$-axis; then
        up $4$.}

  \item \textbf{From the graph.} Write the equation of the function below in the form
        $y=a\sqrt{x-h}+k$.
        \begin{center}
        \begin{tikzpicture}[scale=0.42]
          \lgrid{-2}{13}{-4}{6}
          \begin{scope}
            \clip (-2,-4) rectangle (13,6);
            \draw[very thick, forest, domain=1:13, samples=100, smooth] plot (\x,{2*sqrt((\x)-1)-2});
          \end{scope}
          \fill[forest] (1,-2) circle (4pt) node[below left, font=\scriptsize]{$(1,-2)$};
          \fill[charcoal] (5,2) circle (2.6pt) node[above left, font=\scriptsize]{$(5,2)$};
          \fill[charcoal] (10,4) circle (2.6pt);
        \end{tikzpicture}
        \end{center}
        \vspace{-6pt}
        Anchor: \ans{$(1,-2)$} \; $\Rightarrow$ \; $h=\mans{1}$, $k=\mans{-2}$. \quad
        Then substitute $(5,2)$ to find $a=\mans{2}$.

        \vspace{4pt}
        \textbf{Equation:} $y=$ \ans{$2\sqrt{x-1}-2$}

        \vspace{2pt}
        {\small\emph{Work:} $2=a\sqrt{5-1}-2 \Rightarrow 2=2a-2 \Rightarrow a=2$. \;
        Check $(10,4)$: $2\sqrt{9}-2=4$. \checkmark}

  \item \textbf{SOL-style multiple choice.} What is the anchor point (endpoint) of the graph of
        $y=3\sqrt{2x-10}+1$?
        \begin{enumerate}[label=\Alph*., leftmargin=1.9em, itemsep=1pt]
          \item $(10,1)$
          \item $(5,1)$ \quad \textcolor{keyred}{\textbf{$\leftarrow$ correct}}
          \item $(-5,1)$
          \item $(5,-1)$
        \end{enumerate}
        \vspace{2pt}
        \begin{itemize}[leftmargin=1.4em, nosep]
          \item \textbf{B} is correct: factor first. $\sqrt{2x-10}=\sqrt{2(x-5)}=\sqrt2\,\sqrt{x-5}$,
                so the anchor is $(5,1)$. (Or the old way: $2x-10\ge0 \Rightarrow x\ge5$.)
          \item \textbf{A} read the $10$ straight out of the radicand without factoring out the
                $2$ --- the error Section 5 exists to prevent.
          \item \textbf{C} factored but kept the constant's sign, turning $x-5$ into a shift
                \emph{left}.
          \item \textbf{D} flipped the sign of $k$ as well. Only the \emph{inside} number reverses.
        \end{itemize}
\end{enumerate}

\begin{teachernote}
\small Graded for completion. Four piles, and they decide how 5.5 opens.
\begin{itemize}[leftmargin=1.3em, nosep]
  \item \textbf{Got it} --- 1(a)--(e) right, item 2 $=2\sqrt{x-1}-2$, item 3 $=$ B.
  \item \textbf{Sign-of-$h$ slips} --- 1(a) answered $(2,4)$, or item 2 answered $\sqrt{x+1}$.
        Cheapest fix: set the radicand to zero and solve --- that \emph{is} $h$, and it never lies.
  \item \textbf{Range droppers} --- 1(c) answered $[4,\infty)$: found the anchor, stopped reading.
        Most of this pile got 1(d) right; showing them their own inconsistency does the teaching.
  \item \textbf{Did not factor} --- item 3 $=$ A. The pile that matters: a \emph{procedural} gap that
        resurfaces in 5.5, where $\sqrt{2x-10}=4$ must be handled correctly.
\end{itemize}
\textbf{Item 1(d) is the conceptual item.} A student who names the anchor but cannot say \emph{which
symbol} made it a maximum has memorized a picture, not a mechanism.
\end{teachernote}

\end{document}
